\section{What a budget buys}
\label{sec:agents}

Define the \emph{free compression ratio} of a flat record,
$C(n) = 2^n / \Kflat(n)$. Dividing the bounds of Section~\ref{sec:threshold}
through by $2^n$ gives
$2^{\lfloor \log_2(n/2+1)\rfloor} \le C(n) \le (n+2)/2$. Both ends are
$\Theta(n)$, so the safe zone widens as the team grows---but only linearly,
against an exponentially large baseline.

\begin{table}[t]
\centering
\caption{The same numbers under two normalisations, uniform prior. $(n+2)/2$ is
an upper bound on $C(n)$, not its value; the gap is real at $n=7$.}
\label{tab:ratio}
\small
\begin{tabular}{@{}lrrrrrr@{}}
\toprule
$n$ & 3 & 5 & 7 & 9 & 20 & 50 \\
\midrule
$C(n)$        & 2.00 & 3.20 & 4.00 & 4.00--5.28 & 8--11 & 16--26 \\
$(n+2)/2$     & 2.50 & 3.50 & 4.50 & 5.50 & 11 & 26 \\
free bits     & 1.00 & 1.68 & 2.00 & 2.00--2.40 & 3.00--3.46 & 4.00--4.70 \\
as \% of $n$  & 33.3 & 33.6 & 28.6 & 22--27 & 15.0--17.3 & 8.0--9.4 \\
\bottomrule
\end{tabular}
\end{table}

Both normalisations are correct and they point opposite ways, which is why we
give both. As a \emph{state-space ratio} the safe zone widens: one may compress
$3.2\times$ for free at $n=5$, about $26\times$ at $n=50$. In \emph{bits}---the
unit a record is actually sized in---the free allowance is
$\log_2 C(n) \approx \log_2(n/2)$ out of $n$, so the \emph{fraction} of per-agent
resolution one may discard for free \emph{shrinks} with team size. Linear growth
of a ratio against an exponential baseline is logarithmic growth in bits.

Concretely: at $n = 20$, flatness requires $K$ between $95{,}326$ and
$131{,}072$ labels---retaining $16.5$--$17.0$ of the $20$ bits, i.e.\ $83$--$85\%$
of per-agent resolution. At $n = 50$ it requires retaining $45.3$--$46.0$ of the
$50$ bits, i.e.\ $91$--$92\%$. Two design consequences follow immediately, and
Appendix~\ref{app:design} states the rest. Retaining one effort bit per agent is
flat with margin $1$, so the tension simply does not arise; everything here is
about what happens when that is unaffordable. And an aggregate such as ``how many
agents succeeded'' carries $\log_2(n+1)$ bits---$5.7$ at $n = 50$, about $11\%$
of per-agent resolution---an order of magnitude below the threshold, so by
Corollary~\ref{cor:edge} any record mapping a profile and all $n$ of its
unilateral neighbours to one label admits no positive margin at any price.

Bittensor's concealment interval~\cite{bittensorcommitreveal}, tuned to roughly
$1800$ blocks, is the clearest deployed instance of an operator selecting a point
on this curve: more concealment, more incentive margin, slower evaluation. Its
binding constraint is temporal rather than a size bound, so it evidences the
\emph{consequence} of settling from a merged record while the context-window
argument evidences the \emph{mechanism}.

\subsection{What production agent frameworks log by default}
\label{sec:framework-audit}

The theory bounds what a record must retain; whether deployed systems clear that
bar is a separate, checkable question. We audited seven production multi-agent
frameworks on one criterion: does the \emph{default} structured log let a
principal recover, after the fact, which agent did what---not whether custom
instrumentation could, but what a settlement built on the log as shipped would
see. For AutoGen we did not stop at the documentation: we installed
\texttt{autogen-core}~0.7.5 and drove a real three-agent episode through its own
\texttt{SingleThreadedAgentRuntime}, capturing every event the framework's own
logger emitted (\texttt{experiments/production\_probe/} reproduces the run).

\begin{table}[t]
\centering
\caption{Default per-agent attribution in seven production agent frameworks,
audited 2026-09-14 and, for AutoGen, corrected 2026-09-19 after running the
framework directly.}
\label{tab:frameworks}
\small
\begin{tabular}{@{}lll@{}}
\toprule
Framework & Class & Effort recoverable? \\
\midrule
LangSmith / LangGraph~\cite{langsmithdocs} & SAFE & per-step, per-step cost \\
OpenAI Agents SDK~\cite{openaiagentssdk} & SAFE & per agent run \\
Claude Agent SDK & SAFE & per subagent, own file \\
Microsoft Semantic Kernel~\cite{semantickerneltokens} & SAFE & per-agent token/cost \\
Google Agent Development Kit~\cite{googleadktraces} & SAFE (leaning) & per-agent spans \\
AutoGen~\cite{autogenlogging} & SAFE & per agent, ambient context \\
CrewAI~\cite{crewaitracing,crewaiovercountissue} & MIXED & outputs yes, cost no \\
\bottomrule
\end{tabular}
\end{table}

Six of the seven default to per-agent attribution. That number is a correction,
not a summary: an earlier documentation-based pass classified AutoGen MIXED
because \texttt{LLMCallEvent} and \texttt{ToolCallEvent} carry no agent
identifier in their public constructor signature. Running the framework shows
both populate an \texttt{agent\_id} internally from
\texttt{MessageHandlerContext.agent\_id()} whenever the event is logged inside a
message handler---the normal case---and that field is undocumented. In our run,
$120$ of $120$ \texttt{LLMCallEvent}s and $120$ of $120$
\texttt{ToolCallEvent}s across $40$ episodes carried a working per-agent
identifier. We leave the correction visible because it is itself the point:
auditing logging behaviour from documentation alone is not reliable, which is
grounds for scepticism about any claim in this section that has not been checked
against a live run, including ours.

\paragraph{This cuts against the easy version of our own motivation.} If the
industry default were aggregation, this paper's premise would be free. It is
not: the aggregation we found is a convenience default rather than a binding
budget, recoverable with modest instrumentation, and we do not lean on it. What
the audit establishes is narrower---the failure mode exists at the schema level
in a widely used framework, on the axis the theory predicts, and is invisible to
anyone who has not looked at the field level.
